\section{Kiến thức chuẩn bị}
\label{sec:kien-thuc-chuan-bi}

Chương Ranking nghiên cứu các bài toán học máy trong đó đầu ra của mô hình là một thứ tự xếp hạng giữa các đối tượng. Khác với phân loại, ranking không chỉ hỏi một mẫu thuộc lớp nào, mà quan tâm đối tượng nào nên đứng trước đối tượng nào. Bài toán này xuất hiện trong công cụ tìm kiếm, hệ thống gợi ý, truy xuất thông tin, xếp hạng sản phẩm, xếp hạng phim và phát hiện gian lận.

\subsection{Bài toán xếp hạng}

Gọi \(\mathcal{X}\) là không gian đầu vào. Với một tập hữu hạn gồm \(n\) đối tượng:
\[
    X=\{x_1,x_2,\ldots,x_n\}\subseteq \mathcal{X},
\]
mục tiêu của ranking là tìm một thứ tự sao cho các đối tượng quan trọng hơn hoặc liên quan hơn được đặt ở vị trí cao hơn. Nếu \(x_i\) nên được xếp trước \(x_j\), ta viết:
\[
    x_i \succ x_j.
\]

Một ranking đầy đủ có thể được biểu diễn bởi một hoán vị \(\pi\) của tập chỉ số \(\{1,2,\ldots,n\}\):
\[
    x_{\pi(1)} \succ x_{\pi(2)} \succ \cdots \succ x_{\pi(n)}.
\]

\subsection{Hàm điểm và ranking dựa trên điểm số}

Một cách phổ biến để tạo ranking là học một hàm điểm:
\[
    h:\mathcal{X}\rightarrow\mathbb{R}.
\]
Hàm \(h\) gán cho mỗi đối tượng một điểm số thực. Nếu \(h(x_i)>h(x_j)\), mô hình xếp \(x_i\) cao hơn \(x_j\). Do đó, ranking được tạo bằng cách sắp xếp các đối tượng theo điểm số giảm dần:
\[
    h(x_{\pi(1)})\geq h(x_{\pi(2)})\geq \cdots \geq h(x_{\pi(n)}).
\]

Trong nhiều mô hình, mỗi đối tượng được biểu diễn bởi vector đặc trưng \(x_i\in\mathbb{R}^d\). Một lớp hàm điểm đơn giản và quan trọng là hàm tuyến tính:
\[
    h_w(x)=w^\top x,
\]
trong đó \(w\in\mathbb{R}^d\) là vector tham số cần học.

\begin{figure}[H]
    \centering
    \begin{tikzpicture}[>=Stealth]
        \draw[->, thick] (0,0) -- (10.5,0) node[right] {điểm \(h(x)\)};
        \foreach \x/\label/\score in {1.2/\(x_4\)/0.4, 3.0/\(x_2\)/1.1, 5.2/\(x_5\)/1.8, 7.4/\(x_1\)/2.6, 9.0/\(x_3\)/3.1} {
            \draw[fill=gray!30] (\x,0) circle (3pt);
            \node[above] at (\x,0.05) {\label};
            \node[below] at (\x,-0.05) {\(\score\)};
        }
        \draw[->, thick, dashed] (9.2,0.8) -- (1.0,0.8);
        \node[above] at (5.1,0.8) {sắp xếp giảm dần theo điểm số};
    \end{tikzpicture}
    \caption{Score-based ranking tạo thứ tự bằng cách sắp xếp các đối tượng theo điểm số \(h(x)\).}
    \label{fig:section2-score-axis-ranking-short}
\end{figure}

\subsection{Quan hệ ưu tiên theo từng cặp}

Dữ liệu huấn luyện trong ranking thường được cho dưới dạng các cặp đối tượng:
\[
    (x_i,x'_i,y_i)\in
    \mathcal{X}\times\mathcal{X}\times\{-1,0,+1\}.
\]
Nhãn \(y_i\) biểu diễn quan hệ ưu tiên giữa \(x_i\) và \(x'_i\):
\[
    y_i =
    \begin{cases}
        +1, & \text{nếu } x'_i \text{ được ưu tiên hơn } x_i,\\
        -1, & \text{nếu } x_i \text{ được ưu tiên hơn } x'_i,\\
        0,  & \text{nếu hai đối tượng ngang nhau hoặc không có thông tin.}
    \end{cases}
\]

Với quy ước này, một cặp được xếp đúng khi:
\begin{equation}
    y_i\big(h(x'_i)-h(x_i)\big)>0.
    \label{eq:correct-pairwise-ranking}
\end{equation}
Biểu thức trên kiểm tra xem dấu của hiệu điểm giữa hai đối tượng có phù hợp với nhãn ưu tiên hay không.

\subsection{Lỗi xếp hạng theo từng cặp}

Chất lượng của hàm điểm \(h\) có thể được đo bằng xác suất xếp sai một cặp đối tượng. Gọi \(D\) là phân phối trên \(\mathcal{X}\times\mathcal{X}\) và \(f\) là hàm ưu tiên mục tiêu. Lỗi tổng quát được định nghĩa là:
\begin{equation}
    R(h)
    =
    \Pr_{(x,x')\sim D}
    \left[
        f(x,x')\neq 0
        \ \wedge\
        f(x,x')\big(h(x')-h(x)\big)\leq 0
    \right].
    \label{eq:generalization-pairwise-error}
\end{equation}

Trong thực tế, ta chỉ có mẫu huấn luyện:
\[
    S =
    \big((x_1,x'_1,y_1),\ldots,(x_m,x'_m,y_m)\big).
\]
Lỗi thực nghiệm trên \(S\) là:
\begin{equation}
    \widehat{R}(h)
    =
    \frac{1}{m}
    \sum_{i=1}^{m}
    \mathbf{1}
    \left[
        y_i\neq 0
        \ \wedge\
        y_i\big(h(x'_i)-h(x_i)\big)\leq 0
    \right].
    \label{eq:empirical-pairwise-error}
\end{equation}
Đây là tỉ lệ các cặp huấn luyện bị mô hình xếp sai hoặc không phân biệt được đúng chiều.

\subsection{Margin trong ranking}

Tương tự classification, ranking cũng có thể được phân tích thông qua margin. Với một cặp \((x_i,x'_i,y_i)\), margin theo cặp là:
\begin{equation}
    M_i(h)
    =
    y_i\big(h(x'_i)-h(x_i)\big).
    \label{eq:pairwise-margin}
\end{equation}
Nếu \(M_i(h)>0\), cặp được xếp đúng. Nếu \(M_i(h)\leq 0\), cặp bị xếp sai hoặc hoà điểm. Với một ngưỡng \(\rho>0\), lỗi margin thực nghiệm được viết là:
\begin{equation}
    \widehat{R}_{\rho}(h)
    =
    \frac{1}{m}
    \sum_{i=1}^{m}
    \Phi_{\rho}
    \left(
        y_i\big(h(x'_i)-h(x_i)\big)
    \right),
    \label{eq:empirical-margin-loss}
\end{equation}
trong đó \(\Phi_{\rho}\) là hàm mất mát margin. Trực giác là mô hình không chỉ cần xếp đúng, mà còn cần tạo ra khoảng cách điểm đủ lớn giữa hai đối tượng.

\subsection{Score-based ranking và preference-based ranking}

Có hai thiết lập chính trong ranking. Trong \textit{score-based ranking}, mô hình học một hàm điểm:
\[
    h:\mathcal{X}\rightarrow\mathbb{R},
\]
rồi sinh ranking bằng cách sắp xếp các đối tượng theo điểm số.

Trong \textit{preference-based ranking}, mô hình học trực tiếp một hàm ưu tiên:
\[
    h:\mathcal{X}\times\mathcal{X}\rightarrow[0,1].
\]
Giá trị \(h(u,v)\) càng lớn thì mô hình càng tin rằng \(u\) nên được xếp trước \(v\). Có thể tóm tắt sự khác nhau như sau:
\[
    \text{score-based: } x\mapsto h(x),
    \qquad
    \text{preference-based: } (u,v)\mapsto h(u,v).
\]

\subsection{Bipartite ranking và AUC}

Một trường hợp quan trọng của ranking là \textit{bipartite ranking}, trong đó không gian đầu vào được chia thành hai lớp:
\[
    \mathcal{X}=\mathcal{X}_{+}\cup\mathcal{X}_{-},
    \qquad
    \mathcal{X}_{+}\cap\mathcal{X}_{-}=\varnothing.
\]
Mục tiêu là học một hàm điểm \(h\) sao cho các điểm positive được xếp cao hơn các điểm negative:
\[
    h(x^{+})>h(x^{-}),
    \qquad
    x^{+}\in\mathcal{X}_{+},\ x^{-}\in\mathcal{X}_{-}.
\]

Nếu có \(m\) positive points và \(n\) negative points, lỗi thực nghiệm trong bipartite ranking là:
\begin{equation}
    \widehat{R}(h)
    =
    \frac{1}{mn}
    \sum_{i=1}^{m}
    \sum_{j=1}^{n}
    \mathbf{1}
    \left[
        h(x^{+}_i)<h(x^{-}_j)
    \right].
    \label{eq:bipartite-empirical-error}
\end{equation}

AUC có thể được hiểu là tỉ lệ các cặp positive-negative được xếp đúng:
\begin{equation}
    \mathrm{AUC}(h)
    =
    \frac{1}{mn}
    \sum_{i=1}^{m}
    \sum_{j=1}^{n}
    \mathbf{1}
    \left[
        h(x^{+}_i)\geq h(x^{-}_j)
    \right].
    \label{eq:auc-definition}
\end{equation}
Do đó, nếu bỏ qua trường hợp hoà điểm hoặc xử lý hoà điểm theo một quy ước cố định, ta có:
\begin{equation}
    \mathrm{AUC}(h)=1-\widehat{R}(h).
    \label{eq:auc-error-relation}
\end{equation}